% Zone „Das kennst du schon“ – aus bank/hypergeometrische-verteilung/zone.jsonl (zusammenbau v0.1)
\einheitenkopf[zone][Kennst du schon]{Das kennst du schon}
\setcounter{aufgabe}{0}

%% TODO Titel: Ich-kann-Satz fehlt in der Bank (Platzhalter: Kettenname) – Nr. 1
%% TODO Anweisung: ein Satz über der ersten Teilaufgabe fehlt in der Bank – Nr. 1
\begin{aufgabe}{Binomialkoeffizient (n über k) als Anzahl der Auswahlen, Rechnertaste und kleine Werte im Kopf}
\begin{teile}
\teil Wie viele Möglichkeiten gibt es, aus $5$ Personen $2$ auszuwählen? Das ist $(5 \text{ über } 2)$.
\teil Berechne $(10 \text{ über } 3)$ mit dem Taschenrechner.
\end{teile}
\end{aufgabe}

%% TODO Titel: Ich-kann-Satz fehlt in der Bank (Platzhalter: Kettenname) – Nr. 2
%% TODO Anweisung: ein Satz über der ersten Teilaufgabe fehlt in der Bank – Nr. 2
\begin{aufgabe}{Pfadregeln ohne Zurücklegen (Bruchkette mit schrumpfendem Nenner) und das Gegenereignis}
\begin{teile}
\teil In einer Urne liegen $3$ rote und $2$ blaue Kugeln. Es werden $2$ Kugeln ohne Zurücklegen gezogen. Berechne die Wahrscheinlichkeit, dass beide rot sind.
\teil In einem Beutel sind $4$ grüne und $3$ gelbe Chips. Es werden $3$ Chips ohne Zurücklegen gezogen. Berechne die Wahrscheinlichkeit, dass alle grün sind.
\end{teile}
\end{aufgabe}

%% TODO Titel: Ich-kann-Satz fehlt in der Bank (Platzhalter: Kettenname) – Nr. 3
%% TODO Anweisung: ein Satz über der ersten Teilaufgabe fehlt in der Bank – Nr. 3
\begin{aufgabe}{Das Binomialmodell und seine Bedingungen (zur Abgrenzung: wann es ungeeignet ist)}
\begin{teile}
\teil Ein Würfel wird zehnmal geworfen, $X$ zählt die Sechsen. Ist $X$ binomialverteilt? \\ \kreuz{ja} \\ \kreuz{nein}
\teil In einer Stadt mit $50\,000$ Einwohnern fahren $30\,\%$ mit dem Rad zur Arbeit. $20$ Einwohner werden zufällig befragt, $X$ zählt die Radfahrer. Ist das Binomialmodell als Näherung geeignet? \\ \kreuz{ja} \\ \kreuz{nein}
\end{teile}
\end{aufgabe}

%% TODO Titel: Ich-kann-Satz fehlt in der Bank (Platzhalter: Kettenname) – Nr. 4
%% TODO Anweisung: ein Satz über der ersten Teilaufgabe fehlt in der Bank – Nr. 4
\begin{aufgabe}{Brüche multiplizieren und kürzen}
\begin{teile}
\teil Berechne und kürze: $\frac{2}{3} \cdot \frac{3}{4}$
\teil Berechne und kürze: $\frac{5}{8} \cdot \frac{4}{7} \cdot \frac{3}{6}$
\end{teile}
\end{aufgabe}

%% TODO Titel: Ich-kann-Satz fehlt in der Bank (Platzhalter: Kettenname) – Nr. 5
%% TODO Anweisung: ein Satz über der ersten Teilaufgabe fehlt in der Bank – Nr. 5
\begin{aufgabe}{Pfadregeln ohne Zurücklegen (Bruchkette mit schrumpfendem Nenner) und das Gegenereignis – Fehler finden}
\begin{teile}
\teil Finde den Fehler und rechne richtig. In einer Urne liegen $5$ weiße und $3$ schwarze Kugeln; es werden $2$ ohne Zurücklegen gezogen. Tim rechnet: $P(\text{beide schwarz}) = \frac{3}{8} \cdot \frac{3}{8} = \frac{9}{64}$.
\end{teile}
\end{aufgabe}

%% TODO Titel: Ich-kann-Satz fehlt in der Bank (Platzhalter: Kettenname) – Nr. 6
%% TODO Anweisung: ein Satz über der ersten Teilaufgabe fehlt in der Bank – Nr. 6
\begin{aufgabe}{Pfadregeln ohne Zurücklegen (Bruchkette mit schrumpfendem Nenner) und das Gegenereignis – selbst rechnen}
\begin{teile}
\teil In einer Urne liegen $6$ rote und $2$ gelbe Kugeln; es werden $2$ ohne Zurücklegen gezogen. Berechne die Wahrscheinlichkeit, dass beide gelb sind.
\end{teile}
\end{aufgabe}
